\chapter*{คำนำ}
\addcontentsline{toc}{chapter}{คำนำ}

การตรวจวัดและประเมินคุณสมบัติความอุดมสมบูรณ์ของดินในพื้นที่การเกษตรแม่นยำ เช่น สวนทุเรียนและพืชเศรษฐกิจในเขตภาคตะวันออกของประเทศไทย ถือเป็นหัวใจสำคัญที่มีผลต่อผลผลิตและคุณภาพทางการเกษตรโดยตรง ในอดีตการวิเคราะห์ดินจำเป็นต้องอาศัยการเก็บตัวอย่างส่งห้องปฏิบัติการทางเคมีซึ่งมีค่าใช้จ่ายสูงและใช้ระยะเวลานาน หรือการใช้หัววัดอิเล็กทรอนิกส์ภาคสนามแบบโลหะแทงสัมผัส (Metallic Soil Probe 7-in-1) เพียงอย่างเดียว ซึ่งแม้จะให้ผลการวัดได้ทันที แต่หัววัดภาคสนามยังคงมีข้อจำกัดด้านความแปรปรวนจากแรงกด ความชื้นสัมพัทธ์ในเนื้อดิน และการสึกกร่อนของผิวสัมผัสโลหะเมื่อใช้งานต่อเนื่อง

หนังสือคู่มือวิชาการและวิจัยฉบับนี้ จัดทำขึ้นเพื่อเป็นคู่มือภาคปฏิบัติการและหลักการเชิงทฤษฎีขั้นสูง ในการพัฒนาและประยุกต์ใช้เทคโนโลยีคอมพิวเตอร์วิทัศน์บนสมาร์ทโฟน (Smartphone Computer Vision) ผสานเข้ากับโมเดลโครงข่ายประสาทเทียมเชิงลึกแบบหลากมิติ (Multimodal Deep Neural Network) เพื่อยกระดับการตรวจวิเคราะห์ดินไปสู่มิติใหม่ โดยผสานรวมคุณลักษณะทางแสงและสีของหน้าดินในหลายปริภูมิสี ได้แก่ RGB, HSV, CIE $L^*a^*b^*$ และ YCbCr เข้ากับข้อมูลกายภาพที่ตรวจวัดได้จากเซนเซอร์ภาคสนามแบบเรียลไทม์

เนื้อหาภายในคู่มือเล่มนี้ถูกแบ่งออกเป็น 6 บทหลัก เริ่มต้นตั้งแต่พื้นฐานฟิสิกส์เชิงแสงและสรีรวิทยาของสีดิน การแปลงและคำนวณคุณลักษณะสีตามมาตรฐานสากล สถาปัตยกรรมการเลือกพื้นที่วิเคราะห์แบบสัมผัสหน้าจอ (Interactive Touch ROI Geometry) ทั้งรูปทรงสี่เหลี่ยม วงกลมกล้องจุลทรรศน์ และรูปทรงหลายเหลี่ยมอิสระ การออกแบบสถาปัตยกรรมโครงข่ายประสาทเทียม 16 มิติเพื่อทำนายค่าอินทรียวัตถุ (SOM) ความเป็นกรดด่างเชิงแสง (Optical pH) ไนโตรเจน ฟอสฟอรัส โพแทสเซียม และการตรวจสอบความสอดคล้อง ($\Delta\text{pH}$) ตลอดจนวิศวกรรมการประมวลผลสตรีมภาพสดแบบ Zero-Latency และการสอบเทียบความเที่ยงตรงภาคสนามจริง

ผู้เขียนหวังเป็นอย่างยิ่งว่า คู่มือเล่มนี้จะเป็นประโยชน์สูงสุดต่อนักวิจัย วิศวกรเกษตรดิจิทัล นักศึกษาระดับบัณฑิตศึกษา และผู้สนใจทุกท่าน ในการนำองค์ความรู้ไปต่อยอดและประยุกต์ใช้เพื่อการพัฒนาเทคโนโลยีการเกษตรอัจฉริยะอย่างยั่งยืนสืบไป

\vspace{1.5cm}

\begin{flushright}
\textbf{ผู้ช่วยศาสตราจารย์ ดร.ชีวะ ทัศนา}\\[0.2cm]
กลุ่มวิจัยฟิสิกส์เกษตรและนวัตกรรมเซนเซอร์อัจฉริยะ\\
คณะวิทยาศาสตร์และเทคโนโลยี มหาวิทยาลัยราชภัฏรำไพพรรณี\\
กันยายน 2569
\end{flushright}
